\chapter{Bẫy Tự Ái Và Sĩ Diện}
\markboth{Bẫy Tự Ái Và Sĩ Diện}{The Trap of Ego and Saving Face}

\section{Bị Góp Ý Là Nghe Như Bị Hạ Thấp}
\begin{truyen}{Bị Góp Ý Là Nghe Như Bị Hạ Thấp}{Hearing Feedback as Personal Diminishment}
\chuhoa{C}{ó người chỉ cần bị nhắc một lỗi nhỏ cũng thấy như mình bị đánh giá cả con người.}
\textit{(Some people hear one small correction and feel as though their whole character is being judged.)}

Vì tự ái nổi lên quá nhanh, họ phản ứng trước khi kịp hiểu điều đang được góp ý.
\textit{(Because ego rises too quickly, they react before they understand what is actually being said.)}

Phản ứng đó giúp giữ sĩ diện trong vài phút nhưng làm chậm việc sửa mình trong nhiều tháng.
\textit{(That reaction may save face for a few minutes while delaying real improvement for months.)}

Một cái tôi quá nhạy khiến người ta khó lớn dù rất thông minh.
\textit{(An overly sensitive ego can keep a person from growing even if they are intelligent.)}

\end{truyen}

\begin{baihoc}
Góp ý đúng không làm em nhỏ đi. Nó chỉ cho em một chỗ để lớn thêm.
\textit{(Accurate feedback does not make you smaller. It simply reveals a place where you can grow.)}
\end{baihoc}

\begin{ghinhoanh}
\textbf{Short English to remember:}

Accurate feedback does not make you smaller.

It simply reveals a place where you can grow.

\end{ghinhoanh}

\begin{tuhocnhanh}
\textbf{Câu hỏi tự học / Self-study questions}

1. Nhân vật đang nghe góp ý như thế nào? \textit{How is the character hearing feedback?}

2. Điều gì đang bị chậm lại? \textit{What is being slowed down?}

3. Em cần học nghe điều khó nghe ở đâu? \textit{Where do you need to learn to hear hard truths better?}
\end{tuhocnhanh}
\ngancach

\section{Sai Rõ Ràng Nhưng Không Muốn Xin Lỗi}
\begin{truyen}{Sai Rõ Ràng Nhưng Không Muốn Xin Lỗi}{Clearly Wrong but Unwilling to Apologize}
\chuhoa{C}{ó người biết mình làm chưa đúng, nhưng lời xin lỗi cứ mắc lại ở cổ vì sợ mất mặt.}
\textit{(Some people know they were wrong, but the apology stays stuck in the throat because they fear losing face.)}

Họ nghĩ lùi một bước sẽ làm mình yếu đi.
\textit{(They think stepping back will make them look weak.)}

Thực ra, nhiều mối quan hệ nứt không phải vì lỗi ban đầu quá lớn, mà vì cái tôi không chịu cúi xuống đúng lúc.
\textit{(In reality, many relationships crack not because the first mistake was too big, but because the ego refused to bend at the right time.)}

Xin lỗi không luôn sửa được hết, nhưng không xin lỗi thường làm vết nứt sâu hơn.
\textit{(An apology does not fix everything, but refusing to apologize often deepens the fracture.)}

\end{truyen}

\begin{baihoc}
Người biết nhận sai sớm thường giữ được nhiều điều hơn người chỉ cố giữ mặt.
\textit{(People who admit wrong early often preserve more than those who only try to save face.)}
\end{baihoc}

\begin{ghinhoanh}
\textbf{Short English to remember:}

People who admit wrong early often preserve more than those who only try to save face.

\end{ghinhoanh}

\begin{tuhocnhanh}
\textbf{Câu hỏi tự học / Self-study questions}

1. Điều gì đang chặn lời xin lỗi? \textit{What is blocking the apology?}

2. Vết nứt sâu hơn bằng cách nào? \textit{How does the fracture get deeper?}

3. Em còn nợ ai một lời nhận sai không? \textit{Do you owe anyone an honest apology?}
\end{tuhocnhanh}
\ngancach

\section{Không Biết Thì Vẫn Cố Tỏ Ra Biết}
\begin{truyen}{Không Biết Thì Vẫn Cố Tỏ Ra Biết}{Pretending to Know Instead of Admitting Ignorance}
\chuhoa{N}{hiều người sợ nói “em chưa biết” vì nghĩ câu đó làm mình thấp đi trước người khác.}
\textit{(Many people are afraid to say “I do not know yet” because they believe it lowers them in front of others.)}

Vì vậy, họ trả lời vòng, nói mơ, hoặc giấu khoảng trống bằng vẻ tự tin.
\textit{(So they answer vaguely, speak around the issue, or cover the gap with confidence.)}

Nhưng cái không biết được thú nhận sớm thường học nhanh hơn cái không biết bị che giấu bằng sĩ diện.
\textit{(But ignorance admitted early learns faster than ignorance hidden behind ego.)}

Người thật sự lớn lên không coi việc chưa biết là nhục. Họ coi đó là điểm bắt đầu.
\textit{(People who truly grow do not treat not knowing as shame. They treat it as a starting point.)}

\end{truyen}

\begin{baihoc}
Nói “em chưa biết” đúng lúc là một biểu hiện của trí tuệ sạch.
\textit{(Saying “I do not know yet” at the right time is a sign of clean intelligence.)}
\end{baihoc}

\begin{ghinhoanh}
\textbf{Short English to remember:}

Saying “I do not know yet” at the right time is a sign of clean intelligence.

\end{ghinhoanh}

\begin{tuhocnhanh}
\textbf{Câu hỏi tự học / Self-study questions}

1. Nhân vật đang che điều gì? \textit{What is the character hiding?}

2. Điều gì học nhanh hơn? \textit{What learns faster?}

3. Em có thể thừa nhận mình chưa biết ở đâu? \textit{Where can you honestly admit you do not know yet?}
\end{tuhocnhanh}
\ngancach

\section{Thua Một Lần Là Không Dám Quay Lại}
\begin{truyen}{Thua Một Lần Là Không Dám Quay Lại}{Failing Once and Never Returning}
\chuhoa{C}{ó người thử một lần, vấp một lần, quê một lần, rồi tự hứa sẽ không bao giờ chạm lại chuyện đó nữa.}
\textit{(Some people try once, fail once, feel embarrassed once, and then promise themselves never to touch that thing again.)}

Họ gọi đó là biết thân biết phận.
\textit{(They call it knowing their place.)}

Nhưng nhiều khi đó chỉ là sĩ diện bị đau nên không dám cho mình cơ hội làm lại.
\textit{(But often it is simply wounded pride refusing to grant itself another attempt.)}

Một lần quê không định nghĩa cả năng lực. Nó chỉ thử xem em có đủ khiêm tốn để quay lại hay không.
\textit{(One embarrassing failure does not define your ability. It only tests whether you are humble enough to return.)}

\end{truyen}

\begin{baihoc}
Người vượt lên thường không ít lần xấu hổ hơn người khác. Họ chỉ không lấy xấu hổ làm điểm dừng.
\textit{(People who rise are not always less embarrassed than others. They simply do not make embarrassment their stopping point.)}
\end{baihoc}

\begin{ghinhoanh}
\textbf{Short English to remember:}

They do not make embarrassment their stopping point.

\end{ghinhoanh}

\begin{tuhocnhanh}
\textbf{Câu hỏi tự học / Self-study questions}

1. Điều gì khiến nhân vật không quay lại? \textit{What keeps the character from returning?}

2. Một lần thua thực ra đang thử điều gì? \textit{What is one failure actually testing?}

3. Em đang dừng lại ở việc nào chỉ vì ngại? \textit{What have you stopped only because of embarrassment?}
\end{tuhocnhanh}
\ngancach

\section{Thắng Lý Lẽ Nhưng Mất Người}
\begin{truyen}{Thắng Lý Lẽ Nhưng Mất Người}{Winning the Argument but Losing the Person}
\chuhoa{C}{ó người trong mỗi cuộc va chạm chỉ chăm chăm làm sao để mình không thua lý.}
\textit{(Some people in every conflict focus only on not losing the argument.)}

Họ nói sắc hơn, nhanh hơn, chặt hơn để giữ phần hơn cho mình.
\textit{(They speak sharper, faster, and harder to preserve the upper hand.)}

Nhưng sau khi thắng lời, họ lại để lại phía sau một khoảng cách khó hàn gắn.
\textit{(But after winning the words, they often leave behind a distance that is hard to repair.)}

Tự ái rất thích chiến thắng tức thì, còn quan hệ bền lại cần một kiểu khiêm nhường biết dừng đúng lúc.
\textit{(Pride loves immediate victory, while lasting relationships need the humility to stop at the right moment.)}

\end{truyen}

\begin{baihoc}
Không phải cuộc tranh luận nào đáng để thắng bằng mọi giá.
\textit{(Not every argument is worth winning at any cost.)}
\end{baihoc}

\begin{ghinhoanh}
\textbf{Short English to remember:}

Not every argument is worth winning at any cost.

\end{ghinhoanh}

\begin{tuhocnhanh}
\textbf{Câu hỏi tự học / Self-study questions}

1. Nhân vật đang muốn giữ điều gì? \textit{What is the character trying to preserve?}

2. Điều gì bị mất sau cuộc thắng này? \textit{What is lost after this victory?}

3. Em có đang cố thắng ở một nơi đáng ra nên dừng lại không? \textit{Are you trying to win where you should be pausing?}
\end{tuhocnhanh}
\ngancach

\section{Không Hỏi Vì Sợ Trông Kém}
\begin{truyen}{Không Hỏi Vì Sợ Trông Kém}{Refusing to Ask Because You Fear Looking Weak}
\chuhoa{C}{ó người không hiểu nhưng vẫn im, vì sợ một câu hỏi sẽ làm lộ ra mình chưa đủ giỏi.}
\textit{(Some people do not understand but stay silent because they fear a question will reveal they are not good enough.)}

Sự im lặng đó giữ được mặt mũi trong vài phút.
\textit{(That silence preserves face for a few minutes.)}

Nhưng đổi lại là một chỗ hổng kiến thức kéo dài và một cơ hội học bị bỏ qua.
\textit{(But the trade-off is a long-lasting knowledge gap and a missed chance to learn.)}

Si diện thường làm người ta giữ được vẻ mạnh trong khoảnh khắc, rồi trả giá bằng sự yếu thật lâu về sau.
\textit{(Pride often preserves an image of strength for a moment and then charges interest through long-term weakness.)}

\end{truyen}

\begin{baihoc}
Một câu hỏi đúng lúc giúp em lớn nhanh hơn rất nhiều so với việc giữ vẻ biết rồi.
\textit{(A timely question helps you grow far more than pretending you already know.)}
\end{baihoc}

\begin{ghinhoanh}
\textbf{Short English to remember:}

A timely question grows you more than preserved appearances.

\end{ghinhoanh}

\begin{tuhocnhanh}
\textbf{Câu hỏi tự học / Self-study questions}

1. Nhân vật đang im vì điều gì? \textit{Why is the character staying silent?}

2. Điều gì bị bỏ lỡ? \textit{What is being missed?}

3. Em cần hỏi điều gì mà đang ngại hỏi? \textit{What do you need to ask but still hesitate to ask?}
\end{tuhocnhanh}
\ngancach

\section{Tự Ái Khi Người Khác Làm Tốt Hơn}
\begin{truyen}{Tự Ái Khi Người Khác Làm Tốt Hơn}{Feeling Offended When Others Perform Better}
\chuhoa{C}{ó người nhìn thấy ai đó làm tốt hơn mình là trong lòng lập tức căng lên như vừa bị chạm vào chỗ yếu.}
\textit{(Some people see others perform better and immediately tense up as if a weakness has been touched.)}

Thay vì học, họ bắt đầu phòng thủ hoặc tìm cách chê lại.
\textit{(Instead of learning, they begin defending themselves or looking for faults in return.)}

Khi cái tôi quá nhạy với việc người khác giỏi, nó sẽ biến cơ hội học hỏi thành một mối đe dọa.
\textit{(When the ego is too sensitive to others' excellence, it turns learning opportunities into threats.)}

Người lớn lên nhanh là người đủ bình tĩnh để không coi thành công của người khác là sự xúc phạm mình.
\textit{(Those who grow quickly are calm enough not to treat others' success as a personal insult.)}

\end{truyen}

\begin{baihoc}
Người khác làm tốt hơn em không làm em nhỏ đi. Nó chỉ cho em thấy còn điều để học.
\textit{(Someone doing better than you does not make you smaller. It only shows there is more to learn.)}
\end{baihoc}

\begin{ghinhoanh}
\textbf{Short English to remember:}

Another person's excellence is not your insult.

\end{ghinhoanh}

\begin{tuhocnhanh}
\textbf{Câu hỏi tự học / Self-study questions}

1. Nhân vật phản ứng thế nào trước người giỏi hơn? \textit{How does the character react to someone better?}

2. Cơ hội nào đang bị biến thành đe dọa? \textit{What opportunity is being turned into a threat?}

3. Em có thể học gì từ người đang làm tốt hơn mình? \textit{What can you learn from someone outperforming you?}
\end{tuhocnhanh}
\ngancach

\section{Cố Giữ Hình Ảnh Mạnh Mẽ Nên Không Dám Nhờ Giúp}
\begin{truyen}{Cố Giữ Hình Ảnh Mạnh Mẽ Nên Không Dám Nhờ Giúp}{Maintaining a Strong Image and Refusing Help}
\chuhoa{C}{ó người lúc nào cũng muốn mình trông ổn, trông cứng, trông tự gánh được mọi thứ.}
\textit{(Some people always want to appear steady, strong, and fully self-sufficient.)}

Vì vậy, ngay cả khi đang quá tải họ cũng không dám nói ra hay nhờ một bàn tay hỗ trợ.
\textit{(So even when overloaded, they still do not dare speak up or ask for support.)}

Hình ảnh mạnh mẽ ấy có thể khiến người khác nể, nhưng cũng khiến chính họ mắc kẹt trong cô đơn.
\textit{(That strong image may impress others, but it can trap the person inside loneliness.)}

Có lúc si diện không làm mình mạnh hơn. Nó chỉ làm mình phải chịu một mình lâu hơn.
\textit{(Sometimes pride does not make you stronger. It only makes you endure alone for longer.)}

\end{truyen}

\begin{baihoc}
Nhờ giúp đúng lúc không làm em yếu. Nó giữ em khỏi gãy ở chỗ không cần gãy.
\textit{(Asking for help at the right time does not make you weak. It keeps you from breaking unnecessarily.)}
\end{baihoc}

\begin{ghinhoanh}
\textbf{Short English to remember:}

Asking for help can keep you from breaking alone.

\end{ghinhoanh}

\begin{tuhocnhanh}
\textbf{Câu hỏi tự học / Self-study questions}

1. Nhân vật đang cố giữ hình ảnh gì? \textit{What image is the character trying to maintain?}

2. Điều gì khiến họ mắc kẹt? \textit{What keeps them trapped?}

3. Em cần nhờ giúp ở việc nào mà vẫn đang ngại? \textit{Where do you need help but still hesitate to ask?}
\end{tuhocnhanh}
\ngancach

\section{Sửa Sai Quá Muộn Vì Cần Giữ Mặt}
\begin{truyen}{Sửa Sai Quá Muộn Vì Cần Giữ Mặt}{Correcting Too Late Because You Need to Save Face}
\chuhoa{C}{ó người biết mình đã chọn sai hoặc làm sai, nhưng cứ kéo dài việc sửa vì không muốn mất mặt trước ai đó.}
\textit{(Some people know they chose wrongly or acted wrongly, yet delay correction because they do not want to lose face.)}

Mỗi ngày chậm sửa là mỗi ngày cái sai ăn thêm vào hậu quả.
\textit{(Each day of delay gives the mistake more time to deepen its consequences.)}

Điều đáng tiếc là phần lớn người ta không đánh giá nặng việc sửa sai sớm bằng việc cố chấp quá lâu.
\textit{(The sad part is that most people judge late stubbornness more harshly than early correction.)}

Si diện thường bảo mình giữ mặt thêm chút nữa, trong khi sự khôn ngoan bảo mình quay đầu ngay.
\textit{(Pride tells you to save face a little longer, while wisdom tells you to turn back now.)}

\end{truyen}

\begin{baihoc}
Sửa sớm thường nhẹ hơn cố giữ đúng điều đã sai.
\textit{(Correcting early is often lighter than defending what is already wrong.)}
\end{baihoc}

\begin{ghinhoanh}
\textbf{Short English to remember:}

Correct early; do not waste life defending what is wrong.

\end{ghinhoanh}

\begin{tuhocnhanh}
\textbf{Câu hỏi tự học / Self-study questions}

1. Vì sao nhân vật chậm sửa? \textit{Why is the character delaying correction?}

2. Cái giá của việc kéo dài là gì? \textit{What is the cost of delay?}

3. Em có cần quay đầu sớm ở việc nào không? \textit{Is there anything you need to turn around on sooner?}
\end{tuhocnhanh}
\ngancach

\section{Dám Nhỏ Lại Để Lớn Lên}
\begin{truyen}{Dám Nhỏ Lại Để Lớn Lên}{Becoming Smaller in Ego to Grow Larger in Life}
\chuhoa{C}{ó người chỉ thật sự tiến bộ khi chấp nhận một giai đoạn trông mình chưa giỏi, chưa biết, chưa hay ho như hình ảnh cũ từng giữ.}
\textit{(Some people truly improve only after accepting a season in which they do not look as skilled, informed, or impressive as before.)}

Họ bắt đầu học lại, hỏi lại, sửa lại từ những điều tưởng đã biết.
\textit{(They begin learning again, asking again, and correcting again from things they thought they already knew.)}

Nhờ chịu nhỏ lại ở cái tôi, họ lớn lên rất nhanh ở năng lực và chiều sâu.
\textit{(By letting the ego become smaller, they grow quickly in ability and depth.)}

Không phải cúi xuống là thua. Nhiều khi đó là tư thế tốt nhất để nhặt lại phần trưởng thành mình đã bỏ quên.
\textit{(Bending down is not always losing. Sometimes it is the best posture for picking up the maturity you left behind.)}

\end{truyen}

\begin{baihoc}
Khiêm nhường không làm em bé lại. Nó làm đường lớn lên mở ra.
\textit{(Humility does not make you smaller. It opens the road to growth.)}
\end{baihoc}

\begin{ghinhoanh}
\textbf{Short English to remember:}

Humility opens the road to growth.

\end{ghinhoanh}

\begin{tuhocnhanh}
\textbf{Câu hỏi tự học / Self-study questions}

1. Nhân vật đã dám bỏ điều gì? \textit{What has the character dared to let go of?}

2. Điều gì lớn lên nhờ sự khiêm nhường? \textit{What grows because of humility?}

3. Em cần nhỏ lại ở điểm nào để lớn lên? \textit{Where do you need to become smaller in ego to grow?}
\end{tuhocnhanh}
